\documentclass[11pt,a4paper]{article}
\usepackage[T1]{fontenc}
\usepackage[utf8]{inputenc}
\usepackage[english]{babel}
\usepackage{amsmath,amssymb}
\usepackage[margin=2cm]{geometry}
\setlength{\parindent}{0pt}
\setlength{\emergencystretch}{3em}
\allowdisplaybreaks

\title{About \texttt{maps.tex}}
\author{}
\date{}
\begin{document}
\maketitle
This note explains what the test data in \texttt{maps.tex} is used for. That file is input for the engine in \texttt{main.c}; it is read by the test \texttt{test\_maps\_and\_relations.c} in this folder.

\begin{enumerate}
\item \textbf{Purpose.} The file holds 54 statements about maps, each to be proved. The statements are data: the engine knows none of them in advance.

\item \textbf{What is being tested.} Whether a given map, such as $n \mapsto 3n$ on $\mathbb{N}_{0}$, is injective or surjective; and how images and preimages behave, as in $f(X \cap Y) = f(X) \cap f(Y)$ for every $X$, $Y$ exactly when $f$ is injective.

\item \textbf{How the statements are written.} Every property is written out as a formula, not named: for instance ``$f$ is injective'' is $\forall a, b\ (f(a) = f(b) \implies a = b)$. The engine must work from these definitions.

\item \textbf{How to run it.} From the top folder of the repository:
\begin{verbatim}
cc -std=c11 -O2 -Wall -Wextra -o /tmp/test_maps_and_relations \
    tests/set_theory/test_maps_and_relations.c && /tmp/test_maps_and_relations
\end{verbatim}

\item \textbf{Expected result.} One \texttt{PASS} line for the file when every statement is proved with a written proof. Today the engine proves 10 of the 54 statements: those about a map between sets of numbers that is given by a single formula, such as $n \mapsto 3n$ or $x \mapsto \frac{x + 1}{x - 1}$. Maps defined by cases, with an absolute value, with pairs as values, or between arbitrary sets are reported as not supported, so the test fails. The file states the goal.
\end{enumerate}
\end{document}
